\chapter{软件架构与数据流}

\section{分层结构}

项目的结构可以概括为“配置与语义层—规划层—实体化层—验证层”。

\begin{figure}[H]
\centering
\begin{tikzpicture}[
  node distance=8mm and 11mm,
  box/.style={draw=TGBlue,rounded corners,fill=TGLightBlue,minimum height=10mm,
    text width=3.35cm,align=center,font=\small},
  arrow/.style={-{Latex[length=2.2mm]},thick,color=TGBlue}
]
  \node[box] (cfg) {JSON 工程配置\\YAML 病例语义};
  \node[box,right=of cfg] (analysis) {病例分析\\导管与牙位语义};
  \node[box,right=of analysis] (plan) {七阶段几何规划\\GenerationContext};
  \node[box,below=of plan] (mesh) {Blender / Trimesh\\实体化与布尔运算};
  \node[box,left=of mesh] (stl) {最终 STL\\过程图与 JSON};
  \node[box,left=of stl] (qa) {独立重建期望\\ValidationResult};
  \draw[arrow] (cfg) -- (analysis);
  \draw[arrow] (analysis) -- (plan);
  \draw[arrow] (plan) -- (mesh);
  \draw[arrow] (mesh) -- (stl);
  \draw[arrow] (stl) -- (qa);
  \draw[arrow,dashed] (cfg) |- (qa);
\end{tikzpicture}
\caption{TwinGuide 的主要数据流。规划结果先以数据类保存，实体化阶段才创建网格。}
\label{fig:architecture}
\end{figure}

\subsection{配置与公共类型}

\projectpath{config.py} 定义 \texttt{CaseConfig} 及全部参数数据类；
\projectpath{models.py} 保存导管、导板标架、切割体和输出类型；
\projectpath{types.py} 定义阶段状态和 \texttt{GenerationContext}。这些类型承担了模块间
契约，使上游结果不会通过隐式全局变量传给下游。

\subsection{纯几何规划}

规划层以点、向量、曲线和文件路径作为主要输出。它负责决定“在哪里切、在哪里连、中心线
怎么走”，尽量不在此阶段执行复杂布尔。这样锚点错误可以在调试图中直接看到，而不会被
体素化或网格修复掩盖。

\subsection{实体化与布尔层}

\projectpath{blender/mesh_builders.py} 把圆柱、窗口和中心线转换为封闭网格；
\projectpath{blender/guide_modeling.py} 组织牙列裁剪、体素融合和复切；
\projectpath{blender/booleans.py} 封装 Blender 多求解器和 manifold3d 精确差集。

\subsection{验证层}

\projectpath{guide_validation.py} 不直接相信生成时的对象，而是读取导出的 STL，重新运行
病例分析、牙位映射、切口规划和连接规划，然后以 BVH 查询、中心线采样和拓扑统计检查结果。
这种“重新建立期望”的方式比仅检查文件存在更有价值，但仍有后文所述的覆盖空白。

\section{七个正式阶段与实际实体化步骤}

\begin{table}[H]
\centering
\small
\begin{tabularx}{\textwidth}{c p{0.22\textwidth} X p{0.16\textwidth}}
  \toprule
  阶段 & 名称 & 主要输出 & 代码成熟度\\
  \midrule
  1 & 导管识别与位姿分析 & 两根当前模式导管、导板局部标架 & stable 1.0\\
  2 & 牙位识别 & FDI 映射、牙弓方向、观察窗语义轴 & experimental 0.4\\
  3 & 窗口规划 & 导孔、操作窗、观察窗 cutter & experimental 0.2\\
  4 & 联建锚点 & 导管 Q/P、导板 A 点及配对 & experimental 0.1\\
  5 & 按压梁选点 & 三锚点、Y 汇合点与强化参数 & experimental 0.1\\
  6 & 结构连接 & 四根主梁、Y 梁和特殊梁中心线 & experimental 0.1\\
  7 & 净距调整 & 手机左右摆动封闭包络 & experimental 0.1\\
  \bottomrule
\end{tabularx}
\caption{\texttt{generation\_process.py} 中声明的七阶段。}
\end{table}

七阶段结束后，\texttt{generate\_guide()} 才调用实体化入口。因此“第 7 阶段完成”不等于
STL 已经生成；它只表示所需包络或几何计划已经准备好。真正的切割、融合、导出和渲染发生
在 \texttt{build\_guide\_from\_links()}。

\section{坐标与单位约定}

全部几何使用世界坐标和毫米。局部标架只用来排序、区分方向和构造参数化几何：

\begin{itemize}
  \item \texttt{TemplateFrame.normal} 指向导板牙合外侧；
  \item \texttt{lateral} 与 \texttt{depth} 描述导板局部横向和深度；
  \item 牙位映射给出 \texttt{e\_occ}、每牙 \texttt{local\_tangent} 和
    \texttt{local\_outward}；
  \item 导管轴在病例分析后统一为从 C 口高端指向闭合孔低端。
\end{itemize}

如果输入网格坐标不一致，或轴符号没有被正确统一，后续每个“上/下、内/外、近/远”的
判断都会同时出错。因此项目在牙位和导管阶段均采用失败关闭，而不是继续猜测。

